\section{实验过程与结果}

\subsection{第一版：Codex 逐步构建的尝试与停滞}

在第一版中，使用Codex采用逐步构建策略。先让AI搭建页面骨架，再加入开屏动画和主页背景。但就在这个阶段，已经开始感到"心有余而力不足"——虽然每个元素都是通过自己的Prompt添加的，但对最终成品没有任何清晰的画面。骨架看起来像一个没有个性的模板。如果继续按此模式进行，只能指望每个模块的Prompt都精准命中模糊想象中的效果——这显然不现实。整个过程在非常早期的阶段就停止了：搭骨架 $\rightarrow$ 加开屏动画 $\rightarrow$ 加背景 $\rightarrow$ 感到缺乏方向 $\rightarrow$ \textbf{停滞}。

\begin{figure}[h]
\centering
\fbox{\rule{0pt}{1.5in} \rule{0.8\linewidth}{0pt}}
\caption{Codex第一版首屏效果：仅骨架+开屏动画+背景，缺乏设计辨识度。}
\label{fig:codex}
\end{figure}

分析失败原因，核心在于\textit{没有设计锚点}。初始Prompt仅是"帮我做一个个人主页"这类模糊描述，未先确定配色、字体或设计方向，每次Prompt在独立运作，没有共同参照。虽然没有走到"模块失控"的阶段，但已能预见到：第一，每个新模块都需要极精准的Prompt才能匹配脑海中的模糊想象；第二，没有统一的设计文档，各模块独立调整的结果必然是风格不统一。

\subsection{第二版：Claude Code 设计优先的完整实现}

基于第一版教训，改用Claude Code采用设计优先策略。大约20\%的对话轮次用于编码前的设计讨论。

\textbf{设计定调：}首先确定了"温润 $\times$ 锐气"的核心调性。温润落在材质上——暖米白纸质感底色、衬线字体、充裕留白；锐气落在工艺上——精准动效、横向翻阅、四种不同风格的数据可视化。后期加入第三个关键词"生命律动"——粒子漂浮、唱片旋转、植物摇摆。设计文档的总结是：\textbf{温润材质 $\times$ 锐利工艺 $\times$ 生命律动}。

\textbf{最终成果：}实现了一个包含六个区域、四个主题数据可视化面板、三个页脚彩蛋和隐藏贪吃蛇游戏的主页，已部署至GitHub Pages。代码总量约3,700行（HTML 470行，CSS 约3,100行，JS 约1,700行），完整对话约200+轮。

\begin{figure}[h]
\centering
\fbox{\rule{0pt}{1.5in} \rule{0.8\linewidth}{0pt}}
\caption{第二版首屏：名字逐字飘落动画+菱形装饰线+粒子漂浮。}
\label{fig:v2-hero}
\end{figure}

\subsection{核心功能详解}

\subsubsection{数据可视化}

关于区的四个兴趣标签分别采用四种完全不同的可视化形式：游戏为像素风卡带插槽面板（配合精灵四帧走动动画），音乐为旋转黑胶唱片+12条随机跳动声波，影视为竖排复古电影票根（锯齿撕边+手写体片名），动漫为CRT小电视（扫描线+弹幕漂移+星星闪烁）。四个面板的设计风格由用户提出关键词（"像素风""波形+唱片""电影票""小电视机"），AI负责将关键词转化为具体的CSS布局和JS动画。

\begin{figure}[h]
\centering
\fbox{\rule{0pt}{1.2in} \rule{0.8\linewidth}{0pt}}
\caption{四个数据可视化面板（从左到右）：像素游戏机、环形唱片、电影票根、CRT小电视。}
\label{fig:viz}
\end{figure}

\subsubsection{横向翻阅（招牌交互）}

思考区是整个主页最核心的交互特色。它使用GSAP ScrollTrigger的pin+scrub机制：访客向下滚动页面时，到达思考区后垂直滚动被"接管"转为水平位移，6张文章卡片逐张从右侧滑入。居中卡片呈现聚光灯效果（scale=1，opacity=1），两侧卡片自动减弱（scale=0.94，opacity=0.55）。翻完所有文章后恢复正常垂直滚动。在移动端（640px以下）降级为正常竖直排列，避免横向翻阅在小屏幕上的体验问题。

\begin{figure}[h]
\centering
\fbox{\rule{0pt}{1.2in} \rule{0.8\linewidth}{0pt}}
\caption{思考区横向翻阅：卡片居中呈现，聚光灯效果（中间亮、两侧暗）。}
\label{fig:thoughts}
\end{figure}

\subsubsection{彩蛋系统}

页脚包含五个按发现难度分层的彩蛋：L1时间问候（自动显示，根据浏览器时间切换晨/昼/暮/夜四条问候语）；L2隐藏歌单（点击音符图标展开播放器，内嵌2首mp3）和彩蛋成就（6个徽章，解锁条件包括时间旅人/音乐鉴赏家/秘密读者/留言达人/游戏高手/深夜访客）；L3隐藏一句话（版权行悬停1.5秒后文字淡化重组为悄悄话）；L4贪吃蛇小游戏（连点页脚波浪线5次触发，Canvas实现，蛇身陶土色，食物苔藓绿，含触屏滑动支持）；L5时间预览钟表（页脚按钮可循环切换6个时段，预览全天候色调效果，用于实验报告截图）。

\subsection{体验反馈驱动的迭代}

第一版实现（v1.0）完成后，在浏览器实际使用中发现了若干仅在纸上无法暴露的问题：
\begin{itemize}
    \item 回到顶部按钮颜色与页脚背景融合，几乎不可见——改为40px圆形按钮（带边框和悬停效果）。
    \item 导航链接直接跳转无平滑过渡，与页面整体"温润"调性不匹配——改为window.scrollTo平滑滚动。
    \item 邮箱图标mailto:协议在未配置邮件客户端的电脑上无效——改为点击复制地址+Toast提示。
    \item QQ音乐图标使用了时钟SVG图案——改为双八分音符音乐符号。
    \item 卡片光照默认持续显示，亮度0.25刺眼——改为仅hover触发，亮度降至0.14。
    \item 手记区卡片只有缩略图，无法放大查看——新增全屏灯箱预览功能。
    \item 兴趣标签无点击提示，访客不知道可交互——在标签行左侧添加永久可见的"点开查看$\rightarrow$"标签。
\end{itemize}
这些问题均通过精准的反馈Prompt——描述"看到的问题+期望的效果"——得到修正。

\subsection{两种策略对比}

\begin{table}[h]
\centering
\caption{两种Vibe Coding策略的定性对比。}
\label{tab:comparison}
\resizebox{\columnwidth}{!}{%
\begin{tabular}{p{0.45\columnwidth}p{0.45\columnwidth}}
\hline
\textbf{逐步构建（Codex）} & \textbf{设计优先（Claude Code）} \\
\hline
搭骨架$\rightarrow$逐个加模块$\rightarrow$逐个调效果 & 讨论设计$\rightarrow$写设计文档$\rightarrow$按文档实现 \\
\hline
无全局设计锚点；风格逐渐漂移 & 设计文档锁定全局方向，始终可控 \\
\hline
每个模块效果依赖单独Prompt的精度 & 设计阶段确认方向；实现时文档作参照 \\
\hline
修改成本随模块数增长 & 文档作为纠错参照，修改始终有据可依 \\
\hline
结果：早期停滞，未产出完整成品 & 结果：完整主页，已部署上线 \\
\hline
\end{tabular}%
}
\end{table}

核心差异在于：逐步构建将用户置于"设计师+调试员"的双重角色——一个没有专业背景的用户很难胜任；设计优先则将用户转变为"设计评审"——只需要判断力（"这好不好看"），不需要专业知识（"我该用什么配色"）。从整个过程中学到的最重要经验是：\textbf{在让AI写代码之前，先让AI帮助做设计。}这20\%的前期设计投入，换来的是方向始终受控、修改始终有据、成品远超预期。
